\begin{center}
    {\LARGE \textbf{2011分析化学}} \\[2ex]
    {\large 时量180分钟 \quad 满分90分}
\end{center}

\vspace{0.5cm}
\noindent\textbf{注意事项:}
\begin{enumerate}[label=\arabic*., parsep=0pt, itemsep=0pt]
    \item 答题前,考生先将自己的姓名、学号、考场号、座位号填写在答题卡上,将条形码准确粘贴在考生信息条形码粘贴区。
    \item 考试结束后,将本试卷与答题纸一并收回。
    \item 回答各个问题时,将答案写在答题纸上,写在本试卷上无效。
\end{enumerate}

\vspace{0.5cm}
\noindent\textbf{一、简答题(42 pts.)}

\begin{enumerate}[label=\arabic*., itemsep=1.5ex]
    \item 欲在$\mathrm{pH}=10.0$的氨性缓冲液中，以$2.0\times10^{-3}\ \mathrm{mol/L}\ \mathrm{EDTA}$滴定同浓度的$\ce{Mg^{2+}}$，下列两种铬黑类的指示剂应选用哪一种？（已知铬黑$\mathrm{A(EBA)}$的$\mathrm{p}K_{a1}=6.2$，$\mathrm{p}K_{a2}=13.0$，$\lg K(\ce{Mg\cdot EBA})=7.2$；铬黑$\mathrm{T(EBT)}$的$\mathrm{p}K_{a2}=6.3$，$\mathrm{p}K_{a3}=11.6$，$\lg K(\ce{Mg\cdot EBT})=7.0$；$\lg K(\ce{MgY})=8.7$；$\mathrm{pH}=10$时，$\lg\alpha_{Y(\mathrm{H})}=0.45$）
    \item $0.001\ \mathrm{mol/L}\ \ce{Ca(NO3)2}$与$0.001\ \mathrm{mol/L}\ \ce{(NH4)2C2O4}$等体积混合，在$\mathrm{pH}=1.0$时，是否有沉淀产生？[$K_{\mathrm{sp}}(\ce{CaC2O4})=10^{-8.64}$；$\mathrm{p}K_{a1}(\ce{H2C2O4})=1.22$，$\mathrm{p}K_{a2}(\ce{H2C2O4})=4.19$]
    \item 于$\mathrm{pH}=10.0$时，以$0.0100\ \mathrm{mol/L}\ \mathrm{EDTA}$溶液滴定$50.00\ \mathrm{mL}$同浓度的某金属离子$\ce{M^{2+}}$；已知在此条件下滴定反应进行完全。当加入$\mathrm{EDTA}$溶液$49.95\ \mathrm{mL}$和$50.05\ \mathrm{mL}$时，化学计量点前后的$\mathrm{pM}$值改变$2$个单位，计算络合物$\ce{MY^{2-}}$的稳定常数$K(\ce{MY})$值。[不考虑$\ce{M^{2+}}$的其它副反应，已知$\mathrm{pH}=10.0$时，$\lg\alpha_{Y(\mathrm{H})}=0.5$]
    \item 采用一般分光光度法，在$525\ \mathrm{nm}$处，以$1\ \mathrm{cm}$比色皿测量$1.00\times10^{-3}\ \mathrm{mol/L}\ \ce{KMnO4}$标准溶液和一$\ce{KMnO4}$试液，其吸光度值分别为$0.700$和$1.00$。今采用示差光度法，以$1.00\times10^{-3}\ \mathrm{mol/L}\ \ce{KMnO4}$标准溶液作参比，问该试液的吸光度为多少？
    \item 要使置信度$90\%$的平均值的置信区间不超过$\pm s$，问至少各应平行测定几次？
    \\
    已知：
    \begin{center}
    \begin{tabular}{l|ccccc}
    $f$ & $1$ & $2$ & $3$ & $4$ & $5$ \\
    \hline
    $t_{0.10}$ & $6.31$ & $2.92$ & $2.35$ & $2.13$ & $2.02$ \\
    \end{tabular}
    \end{center}
    \item 若某物质在两相中的分配比为$17$，今有$50\ \mathrm{mL}$的水溶液，应用$10\ \mathrm{mL}$的萃取剂溶液连续萃取几次，才能使总萃取率达$95\%$以上？
\end{enumerate}

\vspace{0.5cm}
\noindent\textbf{二、计算题(48 pts. total)}

\begin{enumerate}[label=\arabic*., itemsep=1.5ex]
    \item (12 pts.) $0.100\ \mathrm{mol/L}\ \ce{NaOH}$溶液滴定$20.00\ \mathrm{mL}\ 0.0500\ \mathrm{mol/L}\ \ce{H2C2O4}$计算：
    \begin{enumerate}[label=(\arabic*), noitemsep, topsep=0pt, partopsep=0pt]
        \item 化学计量点前$0.1\%$的$\mathrm{pH}$
        \item 化学计量点的$\mathrm{pH}$
        \item 化学计量点后$0.1\%$的$\mathrm{pH}$
    \end{enumerate}
    （$\ce{H2C2O4}$的$\mathrm{p}K_{a1}=1.25$，$\mathrm{p}K_{a2}=4.29$）
    \item (12 pts.) 计算在$[\ce{NH3}]=0.2\ \mathrm{mol/L}$，$[\ce{NH4^+}]=0.1\ \mathrm{mol/L}$的$100\ \mathrm{mL}$溶液中最多能溶解$\ce{Ag2S}$多少克？（$K_{\mathrm{sp}}(\ce{Ag2S})=10^{-48.7}$，$K_\mathrm{b}(\ce{NH3})=10^{-4.74}$，$\ce{H2S}$的$\mathrm{p}K_{a1}=6.88$，$\mathrm{p}K_{a2}=14.15$，$\ce{Ag^+}$与$\ce{NH3}$所生成的络合物的$\beta_1=10^{3.24}$，$\beta_2=10^{7.06}$；$M_r(\ce{Ag2S})=247.8$）
    \item (14 pts.) 用$\mathrm{pH}=5.5$，$0.02\ \mathrm{mol/L}\ \mathrm{EDTA}$滴同浓度$\ce{Zn}$，$\ce{Al}$溶液中的$\ce{Zn}$用$\ce{KF}$掩蔽$\ce{Al}$，终点时$[\ce{F^-}]=10^{-2}\ \mathrm{mol/L}$：$\mathrm{pZn_{sp}}=5.7$，计算终点误差。($\mathrm{pH}=5.5$时，$\lg\alpha_{Y(\mathrm{H})}=5.5$，$\lg K(\ce{ZnY})=16.50$；$\lg K(\ce{AlY})=16.3$；$\ce{Al^{3+}}$和$\ce{F^-}$所生成的络合物的$\beta_1=10^{6.31}$，$\beta_2=10^{11.15}$，$\beta_3=10^{15.00}$，$\beta_4=10^{17.75}$，$\beta_5=10^{19.37}$，$\beta_6=10^{19.84}$)
    \item (10 pts.) 计算在$1\ \mathrm{mol/L}\ \ce{HCl}$介质中$\ce{Fe^{3+} + e^- = Fe^{2+}}$的条件电位。($\ce{Fe^{2+}}$与$\ce{Cl^-}$所生成络合物的$\beta_1=10^{0.4}$，$\ce{Fe^{3+}}$与$\ce{Cl^-}$所生成络合物的$\beta_1=10^{0.6}$，$\beta_2=10^{0.7}$，$\beta_3=10^{-0.7}$，$E^\ominus_{\ce{Fe^{3+}/Fe^{2+}}}=0.771\ \mathrm{V}$)
\end{enumerate}